\documentclass[10pt,a4paper]{amsart}
\usepackage{amsmath,amssymb,amsthm,graphicx}
\usepackage{verbatim}
\newenvironment{reference}{\comment}{\endcomment}
%\newenvironment{reference}{}{}
\newenvironment{slogan}{\comment}{\endcomment}
\newenvironment{history}{\comment}{\endcomment}
% For commutative diagrams we use Xy-pic
\usepackage[all]{xy}

% We use 2cell for 2-commutative diagrams.
\xyoption{2cell}
% Theorem environments.
%
\theoremstyle{plain}
\newtheorem{theorem}[subsection]{Theorem}
\newtheorem{proposition}[subsection]{Proposition}
\newtheorem{lemma}[subsection]{Lemma}

\theoremstyle{definition}
\newtheorem{definition}[subsection]{Definition}
\newtheorem{example}[subsection]{Example}
\newtheorem{exercise}[subsection]{Exercise}
\newtheorem{situation}[subsection]{Situation}

\theoremstyle{remark}
\newtheorem{remark}[subsection]{Remark}
\newtheorem{remarks}[subsection]{Remarks}

\numberwithin{equation}{subsection}

% Macros
%
\def\lim{\mathop{\mathrm{lim}}\nolimits}
\def\colim{\mathop{\mathrm{colim}}\nolimits}
\def\Spec{\mathop{\mathrm{Spec}}}
\def\Hom{\mathop{\mathrm{Hom}}\nolimits}
\def\Ext{\mathop{\mathrm{Ext}}\nolimits}
\def\SheafHom{\mathop{\mathcal{H}\!\mathit{om}}\nolimits}
\def\SheafExt{\mathop{\mathcal{E}\!\mathit{xt}}\nolimits}
\def\Sch{\mathit{Sch}}
\def\Mor{\mathop{\mathrm{Mor}}\nolimits}
\def\Ob{\mathop{\mathrm{Ob}}\nolimits}
\def\Sh{\mathop{\mathit{Sh}}\nolimits}
\def\NL{\mathop{N\!L}\nolimits}
\def\CH{\mathop{\mathrm{CH}}\nolimits}
\def\proetale{{pro\text{-}\acute{e}tale}}
\def\etale{{\acute{e}tale}}
\def\QCoh{\mathit{QCoh}}
\def\Ker{\mathop{\mathrm{Ker}}}
\def\Im{\mathop{\mathrm{Im}}}
\def\Coker{\mathop{\mathrm{Coker}}}
\def\Coim{\mathop{\mathrm{Coim}}}

% Boxtimes
%
\DeclareMathSymbol{\boxtimes}{\mathbin}{AMSa}{"02}

%
% Macros for moduli stacks/spaces
%
\def\QCohstack{\mathcal{QC}\!\mathit{oh}}
\def\Cohstack{\mathcal{C}\!\mathit{oh}}
\def\Spacesstack{\mathcal{S}\!\mathit{paces}}
\def\Quotfunctor{\mathrm{Quot}}
\def\Hilbfunctor{\mathrm{Hilb}}
\def\Curvesstack{\mathcal{C}\!\mathit{urves}}
\def\Polarizedstack{\mathcal{P}\!\mathit{olarized}}
\def\Complexesstack{\mathcal{C}\!\mathit{omplexes}}
% \Pic is the operator that assigns to X its picard group, usage \Pic(X)
% \Picardstack_{X/B} denotes the Picard stack of X over B
% \Picardfunctor_{X/B} denotes the Picard functor of X over B
\def\Pic{\mathop{\mathrm{Pic}}\nolimits}
\def\Picardstack{\mathcal{P}\!\mathit{ic}}
\def\Picardfunctor{\mathrm{Pic}}
\def\Deformationcategory{\mathcal{D}\!\mathit{ef}}
\title{Over a base containing the rationals the de Rham complex of a line bundle is quasi-isomorphic to that of the base, via an explicit homotopy}
\author{The Stacks Project Authors}
\date{stacks-project:derham.tex:local-snapshot-20261012; GNU FDL 1.2 or later}
\begin{document}
\maketitle
\noindent\textit{Editorial note.} All delivered mathematical text is the author's own, copied from the pinned Stacks Project chapter file derham.tex; no statement or proof step was rewritten, reordered or removed.\par\medskip
\noindent\textit{Editorial note.} Preamble provenance: the pinned chapter file opens with \textbackslash{}input\{preamble\} and carries no preamble text of its own, so the project's single Stacks-wide preamble staged at sources/stacks-derham/preamble.tex is excerpted verbatim as three spans, lines 8-12 (verbatim package and the project's own reference/slogan/history display environments), lines 14-18 (Xy-pic setup) and lines 154-217 (theorem declarations and mathematical macros). The wrapper class and the amsmath/amssymb/amsthm/graphicx base packages are supplied by the assembler. No macro was invented or redefined.\par\medskip
\noindent\textit{Editorial note.} Same-chapter pointers are located at their \textbackslash{}label lines inside the pinned chapter file; cross-chapter pointers are located in the project's own tag list staged in the snapshot (sources/stacks-chow/tags.txt), at the verbatim lines carrying their canonical tags.\par\medskip
\noindent\textit{Allowed standard prerequisites (editorial).}\begin{itemize}
\item The situation of the lemma: the bundle, the projection and the de Rham complexes. Same chapter derham.tex, the setup described above the lemma (used at proof lines 1529-1531). It fixes the morphism \$S \textbackslash{}to \textbackslash{}Spec(\textbackslash{}mathbf\{Q\})\$, the line bundle situation and the complexes \$\textbackslash{}Omega\textasciicircum{}\textbackslash{}bullet\_\{X/S\}\$ and \$\textbackslash{}pi\_*\textbackslash{}Omega\textasciicircum{}\textbackslash{}bullet\_\{L/S\}\$ whose comparison is at stake.
\item The affine local computation for a polynomial algebra over a \$\textbackslash{}mathbf\{Q\}\$-algebra. Standard: Kaehler differentials of \$A[t]\$ over \$R\$ and the de Rham complex of a polynomial ring (used at proof lines 1535-1543). It reduces the claim to the assertion that \$\textbackslash{}Omega\textasciicircum{}\textbackslash{}bullet\_\{A/R\} \textbackslash{}to \textbackslash{}Omega\textasciicircum{}\textbackslash{}bullet\_\{A[t]/R\}\$ is a quasi-isomorphism, with the displayed homotopy inverse.
\item Integration of polynomials with rational coefficients. Standard, used at proof lines 1544-1548. The homotopy uses the unique polynomial \$\textbackslash{}int g'\$ with vanishing constant term whose derivative is \$g'\$, which exists because \$R\$ contains \$\textbackslash{}mathbf\{Q\}\$.
\end{itemize}\medskip
\begin{lemma}
\label{lemma-line-bundle-characteristic-zero}
In the situation above, assume there is a morphism $S \to \Spec(\mathbf{Q})$.
Then $\Omega^\bullet_{X/S} \to \pi_*\Omega^\bullet_{L/S}$ is a
quasi-isomorphism and $H_{dR}^*(X/S) = H_{dR}^*(L/S)$.
\end{lemma}

\begin{proof}
Let $R$ be a $\mathbf{Q}$-algebra. Let $A$ be an $R$-algebra.
The affine local statement is that the map
$$
\Omega^\bullet_{A/R} \longrightarrow \Omega^\bullet_{A[t]/R}
$$
is a quasi-isomorphism of complexes of $R$-modules. In fact it is a
homotopy equivalence with homotopy inverse given by the map sending
$g \omega + g' \text{d}t \wedge \omega'$ to $g(0)\omega$ for
$g, g' \in A[t]$ and $\omega, \omega' \in \Omega^\bullet_{A/R}$.
The homotopy sends $g \omega + g' \text{d}t \wedge \omega'$
to $(\int g') \omega'$ where $\int g' \in A[t]$ is the polynomial
with vanishing constant term whose derivative with respect to $t$
is $g'$. Of course, here we use that $R$ contains $\mathbf{Q}$
as $\int t^n = (1/n)t^{n + 1}$.
\end{proof}

\end{document}
